\chapter{Theoretical Fundamentals of Quantum Computing}
\label{c:content}

%%%%%%%%%%%%%%%%%%%%%%%%%%%%%%%%%%%%%%%%%%%%%%%%%%%%%%%%%%%%%%%%%%%%%%%%%%%%%%%%
\section{Basics}
\subsection{Quantum States}
In an $n$-dimensional Hilbert space with a chosen orthonormal basis, a quantum state, denoted by the ket $|\psi\rangle$, can be represented as an $n$-dimensional complex column vector $\vec{\psi}$:
\begin{equation}
	\ket{\psi} \doteq \vec{\psi} = \begin{pmatrix} c_1 \\ c_2 \\ \vdots \\ c_n \end{pmatrix}.
\end{equation}

Its corresponding dual state, denoted by the bra $\langle\psi|$, is represented by the conjugate transpose of this column vector, resulting in an $n$-dimensional row vector $\vec{\psi}^\dagger$:
\begin{equation}
	\langle\psi| \doteq \vec{\psi}^\dagger = \begin{pmatrix} c_1^* & c_2^* & \dots & c_n^* \end{pmatrix}
\end{equation}
\cite[p. 16f.]{Kottmann2025}, \cite[p. 10f.]{sakurai2020modern}.

To ensure consistency with the probability interpretation, quantum states must be normalized to unity. Therefore, for a quantum state $\ket{\psi}$, the sum of the absolute squares of its probability amplitudes must equal one:
\begin{equation}
	\abs{c_1}^2 +\abs{c_2}^2 + \dots + \abs{c_n}^2 = 1 
\end{equation}
\cite[p. 5]{Kottmann2025}.

\subsection{Operators and Observables}
A linear operator acts on a ket from the left to produce another ket vector, or on a bra from the right to produce another bra vector:
\begin{align}
	X\ket{\alpha} &= \ket{\beta}, \\
	\bra{\alpha}X &= \bra{\beta}.
\end{align}

An operator $X$ is defined as the \textbf{null operator} if, for any arbitrary ket $\ket{\alpha}$, it satisfies:
\begin{equation}
	X \ket{\alpha} = 0.
\end{equation}

Operators can be added ($X + Y = Y + X$) and multiplied. While operator multiplication is generally noncommutative ($XY \neq YX$), it is strictly associative:
\begin{equation}
	X(YZ) = (XY)Z = XYZ 
\end{equation} 
\cite[p. 11, 13-15]{sakurai2020modern}.

Because quantum states are represented as vectors in $\mathbb{C}^n$, a linear operator can be represented by a complex $n \times n$ matrix $U$. Given an input ket $\ket{\psi}$, the resulting state $\ket{\phi}$ is calculated via matrix-vector multiplication:
\begin{equation}
	\ket{\phi} = U\ket{\psi}.
\end{equation}

Crucially, physical operations that evolve a quantum state must preserve its probability normalization. Consequently, physical evolution operators are restricted to unitary matrices, ensuring that the inner product—and thus the normalization—remains invariant 
\cite[p. 19]{Kottmann2025}.

An Observable of a system are physical quantities that can be observed within this system. Classical examples are momentum, angular momentum or just the postion or the energy. On a N-Qubit System an observable is defined as a $N \times N$ hermitian matrix.
% ToDo, needs to be finished





\subsection{Unitary Operation and Generators}
To build a quantum circuit, we need a way to express the quantum gates we use later as our operators. Every gate can be written in the following manner:
\begin{equation}
	U(\theta) = e^{-i\frac{\theta}{2}G}
\end{equation}
here $G$ can be any hermitian matrix. Here $G$ is called the Generator.
\cite[p. 34, 39]{Kottmann2025} \cite[p. 1]{kottmann2020feasible}








\subsection{Tensor Products}
% multi-particle system's are constructed in Hilbert space using tensor products (maybe)









